\documentclass[11pt]{article}

\usepackage[spanish]{babel}
\usepackage[utf8]{inputenc}
\usepackage[T1]{fontenc}
\usepackage{geometry}
\usepackage{booktabs}
\usepackage{graphicx}
\usepackage{hyperref}
\usepackage{amsmath}
\usepackage{float}

\geometry{margin=2.5cm}

\title{Proyecto Integrador de Teoría del Riesgo\\
Dashboard de Riesgo Financiero para un Portafolio Internacional}
\author{Nombre del estudiante / grupo}
\date{\today}

\begin{document}

\maketitle

\begin{abstract}
Este trabajo desarrolla un tablero interactivo para el análisis de riesgo financiero de un portafolio compuesto por cinco acciones internacionales: Seven \& i Holdings, Alimentation Couche-Tard, FEMSA, BP y Carrefour. El proyecto integra descarga automática de datos desde APIs, análisis técnico, estudio de rendimientos, modelos ARCH/GARCH, estimación CAPM, medidas VaR/CVaR, optimización de Markowitz, señales automáticas y comparación contra benchmark junto con contexto macroeconómico.
\end{abstract}

\section{Introducción}
La gestión del riesgo financiero requiere integrar herramientas estadísticas, financieras y computacionales para apoyar la toma de decisiones. En este proyecto se construye un dashboard en Streamlit con actualización dinámica de datos y módulos especializados para evaluar diferentes dimensiones del riesgo.

\section{Datos y APIs}
Los precios históricos se obtuvieron mediante Yahoo Finance usando la librería \texttt{yfinance}. Para el contexto macroeconómico se utilizó la API de FRED, incluyendo la serie de tasa libre de riesgo a 3 meses, inflación y tipo de cambio COP/USD.

\section{Activos analizados}
\begin{itemize}
    \item Seven \& i Holdings (\texttt{3382.T})
    \item Alimentation Couche-Tard (\texttt{ATD.TO})
    \item FEMSA (\texttt{FEMSAUBD.MX})
    \item BP (\texttt{BP.L})
    \item Carrefour (\texttt{CA.PA})
\end{itemize}

\section{Metodología}
\subsection{Análisis técnico}
Se calcularon SMA, EMA, RSI, MACD, Bandas de Bollinger y Estocástico.

\subsection{Rendimientos}
Se estimaron rendimientos simples y logarítmicos, junto con estadísticos descriptivos y pruebas de normalidad.

\subsection{Volatilidad condicional}
Se ajustaron modelos ARCH(1), GARCH(1,1) y EGARCH(1,1), comparados mediante AIC y BIC.

\subsection{CAPM}
Se estimó beta para cada activo frente a un benchmark local y se calculó el rendimiento esperado con CAPM.

\subsection{VaR y CVaR}
Para la cuantificación del riesgo extremo del portafolio se definió la variable aleatoria de rendimiento del portafolio en el tiempo \(t\) como
\[
R_{p,t}=\sum_{i=1}^{n} w_i R_{i,t},
\]
donde \(w_i\) representa el peso del activo \(i\) en el portafolio y \(R_{i,t}\) su rendimiento en el periodo \(t\). A partir de esta definición, la variable de pérdida se modeló como
\[
L_t=-R_{p,t},
\]
de modo que valores positivos de \(L_t\) representan pérdidas y valores negativos representan ganancias.

Bajo esta convención, el \textit{Value at Risk} al nivel de confianza \(\alpha\in(0,1)\) se define como el cuantil de orden \(\alpha\) de la distribución de pérdidas:
\[
VaR_{\alpha}(L)=\inf\{l\in\mathbb{R}: P(L\leq l)\geq \alpha\}.
\]
En términos prácticos, \(VaR_{\alpha}\) representa la pérdida máxima esperada que no se excede con probabilidad \(\alpha\) en un horizonte de un periodo. Por ejemplo, un \(VaR_{0.95}=0.025\) indica que, con un 95\% de confianza, la pérdida del portafolio no superaría el 2.5\% en un día.

Como complemento, se empleó el \textit{Conditional Value at Risk} (CVaR), también conocido como \textit{Expected Shortfall}, definido como la pérdida esperada condicionada a que la pérdida exceda el VaR:
\[
CVaR_{\alpha}(L)=E[L \mid L \geq VaR_{\alpha}(L)].
\]
Esta medida resulta especialmente importante en teoría del riesgo porque no solo identifica un umbral crítico, sino que además cuantifica la severidad promedio de las pérdidas extremas más allá de dicho umbral.

Para estimar \(VaR\) y \(CVaR\) se utilizaron tres enfoques complementarios:
\begin{itemize}
    \item \textbf{Método paramétrico:} asume que los rendimientos del portafolio siguen aproximadamente una distribución normal, por lo que el VaR se estima a partir de la media, la desviación estándar y el cuantil de la normal estándar.
    \item \textbf{Método histórico:} no impone una distribución teórica, sino que calcula el cuantil empírico de la muestra observada de pérdidas.
    \item \textbf{Método Monte Carlo:} genera escenarios simulados de rendimientos a partir de una distribución normal multivariada calibrada con la media y la matriz de covarianza muestral de los activos, obteniendo así una distribución simulada de pérdidas del portafolio.
\end{itemize}

La comparación entre estos tres métodos permite evaluar la sensibilidad de las medidas de riesgo frente a diferentes supuestos de modelación. Mientras el enfoque paramétrico depende fuertemente de la hipótesis de normalidad, el histórico recoge directamente la experiencia muestral y el método Monte Carlo permite construir escenarios sintéticos consistentes con la estructura de dependencia estimada entre los activos.

Adicionalmente, se utilizó el test de Kupiec para evaluar la capacidad de cobertura del VaR estimado, contrastando si la frecuencia observada de excedencias es consistente con el nivel de confianza seleccionado.

\subsection{Markowitz}
Se simularon 10{,}000 portafolios para construir el conjunto factible y la frontera eficiente.

\subsection{Señales automáticas}
Se implementaron reglas de compra/venta usando MACD, RSI, Bollinger, cruces de medias y Estocástico.

\subsection{Benchmark y contexto macro}
Se comparó el portafolio contra un benchmark global y se reportaron métricas como Sharpe, Alpha de Jensen, Tracking Error, Information Ratio y máximo drawdown.

\section{Resultados}
En esta sección deben insertarse tablas y figuras del dashboard:
\begin{itemize}
    \item Figura de precios e indicadores.
    \item Tabla descriptiva de rendimientos.
    \item Comparación de modelos GARCH.
    \item Tabla de VaR/CVaR.
    \item Frontera eficiente y pesos óptimos.
    \item Señales por activo.
    \item Desempeño contra benchmark.
\end{itemize}

\section{Conclusiones}
El dashboard permite integrar diferentes enfoques de medición del riesgo en una sola herramienta interactiva. Esto facilita comparar activos, medir exposición sistemática, cuantificar pérdidas extremas y construir portafolios más eficientes.

\section*{Repositorio y ejecución}
\begin{itemize}
    \item Repositorio: completar URL
    \item Ejecución: \texttt{streamlit run app.py}
\end{itemize}

\end{document}
